ソーシャルメディアの登場により、インターネット上では多様な意図を持つ投稿が大量に発信されています。
健全で安全なオンライン空間の実現のためにも、このような投稿に込められた「意図」を解析する技術はますます重要になっています。

私たちは、意図解析技術を題材とする国際コンペティション「SemEval 2020」に参加しました。
最先端の言語モデルにタスク特化型のアーキテクチャを組み合わせ、さらに独自のレシピに基づいて複数のモデルをアンサンブルすることで、高い精度を達成しました。
その結果、\textbf{複数のタスクで1位}を獲得しました。

また、意図を正確に捉えるためには、文の意味構造を解析することも重要です。
私たちは、この構造解析(パーシング)を題材とする「CoNLL 2020 Shared Task」にも参加し、1位を獲得しました。

なお、これらの経験を通じて、アンサンブル学習への関心が深まり、後の研究テーマへとつながりました。
